\subsection{Fasting and Almsgiving}\label{sec:fasting}

St Leo the Great, in a sermon on the December fast, spoke of it as food. In
his Latin, \latin{Semper enim virtuti cibus jejunium fuit}: fasting has always been food for virtue. The English of the Breviary
puts it like this:

\begin{quote}
Fasting hath ever been the bread of strength. From abstinence proceed pure
thoughts, reasonable desires, and healthy counsels. By voluntary mortifications
the flesh dieth to lust, and the soul is renewed in might.\par
{\smallskip\noindent\small ---St Leo the Great, Sermon 13\par}
\end{quote}

What a Christian goes without in Advent feeds something else: the mind that can pray, the will that can choose, and, if the
money saved is given away, the poor. Leo never let the fast stand alone. In
another December sermon he adds its two companions: ``there are three things
which most belong to religious actions, namely prayer, fasting, and
almsgiving''. And he says what each does: ``by prayer we seek to propitiate
God, by fasting we extinguish the lusts of the flesh, by alms we redeem our
sins''.

\subsubsection*{What the law asked, and what it asks}

The law on fasting has changed since the books of 1962, and what bound then is
not what is required now. Here is what the Church required then, and what she
requires now:

\begin{center}
\small
\begin{tabular}{@{}>{\raggedright\arraybackslash}p{0.22\linewidth}>{\raggedright\arraybackslash}p{0.34\linewidth}>{\raggedright\arraybackslash}p{0.34\linewidth}@{}}
\toprule
& Under the Code of 1917, the law of the 1962 books & Now, under the Code of 1983 \\
\midrule
Advent as a season & No fast. & No fast. \\
\addlinespace
Ember Wednesday, Friday and Saturday & Fast and abstinence. & No obligation of their own. \\
\addlinespace
The Vigil, 24 December & Fast and abstinence; from 1959 it might be kept on 23 December instead; no fast when the Vigil fell on a Sunday. & No obligation of its own. \\
\addlinespace
The Fridays of Advent & Abstinence, as on every Friday, unless a holy day of obligation falls on it. & As every Friday of the year, a day of penance. \\
\bottomrule
\end{tabular}
\end{center}

An Ember Friday, or a Vigil that falls on a Friday, is still a Friday. Some
places kept fasts of their own: in the United States the \textit{Manual of
Prayers} of 1889 listed the Fridays of Advent among the fasting days, though
not in some dioceses, such as the Province of New Orleans.

In the United States the bishops wrote in 1966, ``We impose no fast before any
feast-day'', but they left the devout free to prepare for each feast ``by a day
of particular self-denial, penitential prayer and fasting''. That is the place
of fasting in Advent now: not a precept but an offering, made freely, to a
measure that health, age, work and the care of others allow.

\subsubsection*{A fast of one's own}

Keep some fast in Advent all the same. The Ember Days and the Vigil are the
Church's own days for it; a plain
table on those days, or a single full meal, joins the Christian at home to
everyone who has kept them before. The Fridays of Advent can be kept more
strictly than the rest of the year's Fridays. Whatever the measure, Leo's rule
holds: when the body fasts from food, the mind should fast from its vices. A
fast that makes one harsh with one's family has missed the point.

\subsubsection*{Alms}

Leo joins alms to the December fast again and again. In the sermon that speaks of the four fasts of the year he writes:

\begin{quote}
And hence we warn you, beloved, in fatherly affection, to make this winter fast
fruitful to yourselves by bounteous alms, rejoicing that by you the Lord feeds
and clothes His poor, to whom assuredly He could have given the possessions
which He has bestowed on you, had He not in His unspeakable mercy wished to
justify them for their patient labour, and you for your works of love.\par
{\smallskip\noindent\small ---St Leo the Great, Sermon 19\par}
\end{quote}

Elsewhere Leo lists the works that go with the fast: ``Let us defend the widow and serve the
orphan; let us comfort the afflicted and reconcile the estranged; let us take
in the wanderer and succour the oppressed; let us clothe the naked and cherish
the sick.'' And to those who fear to give too much: ``Be stedfast, Christian
giver: give what you may receive, sow what you may reap, scatter what you may
gather.'' Supplications, he says, that are ``winged by such aids mount swiftly
to God's ears''.

Set aside the money that the fast saves, and give it before Christmas. Give
some of it in person. Money sent away does real good, but a gift put into
someone's hand lets the giver see whom the Lord is feeding, and lets the one who
receives it see that someone came. Choose one person you can reach, a neighbour
who lives alone, a family short of money, someone sick or in prison, and bring
the gift yourself before the feast.
